\section*{Step 313 Updated Theorem Audit}

\paragraph{Gaussian magnitude saddle.}
For
\[
T_{k,j}={k\choose j}\zeta^{(j)}(\rho_1)M(G_\star)^{(k-j)}(\rho_1),
\]
the binomial component contributes
\[
\Delta_j^2 \log {k\choose j}
=
\log {k\choose j+1}-2\log {k\choose j}+\log {k\choose j-1}
\approx -\frac{4}{k}
\]
near \(j=k/2\).  Hence the binomial-only Gaussian width is
\[
\sigma_{\rm bin}(k)\approx \sqrt{k}/2.
\]
The numerical decomposition in this step shows that the zeta-derivative
and Mellin-derivative corrections are comparable but do not destroy the
negative quadratic saddle in the certified range.

\paragraph{Phase law.}
The local phase derivative
\[
\alpha(k)\approx \frac12\left(\arg T_{k,j_\ast+1}-\arg T_{k,j_\ast-1}\right)
\]
matches the empirical central phase slope only at a numerical level.
No analytic formula or uniform bound for \(\alpha(k)\) was derived.

\paragraph{Theorem status.}
The Step 310 lower-bound theorem is not upgraded to rigorous status in
this step.  The Gaussian magnitude saddle is structurally supported by
the Stirling/binomial curvature and verified numerically after adding
zeta/M corrections.  The remaining load-bearing gap is a uniform,
analytic control of
\[
\Delta_j \arg\left(\zeta^{(j)}(\rho_1)M(G_\star)^{(k-j)}(\rho_1)\right)
\]
and the discrete stationary-phase remainder.
